%%%%%%%%%%%%%%%%%%%%%%%%%%%%%%%%%%%%%%%%%%%%%%%%%
\section{Introduction}
\label{sec:intro}

Parton distribution functions (PDFs), $f_{i/h}(x,\mu^2)$, are defined as the probability distributions to find a quark, an antiquark, or a gluon ($i=q,\bar{q},g$, respectively) in a fast moving hadron $h$ to carry the hadron's momentum fraction between $x$ and $x+dx$,
probed at the factorization scale $\mu$ \cite{Collins:1989gx}. They are fundamental and important nonperturbative functions in QCD quantifying the relation between a hadron and the quarks and gluons within it, and playing an essential role to connect colliding hadron(s) to short-distance QCD dynamics in high energy scattering processes \cite{Brambilla:2014jmp}. However, calculation of PDFs from the first principle, both analytically or from lattice QCD (LQCD), is a challenge, if it is not impossible, due to the fact that PDFs contain the dynamics at long-distance scales and of nonperturbative in nature, and are defined in terms of time-dependent operators.  Traditionally, PDFs have been extracted from high energy scattering data by QCD global analysis in the framework of QCD factorization \cite{Dulat:2015mca,Martin:2009iq,Ball:2014uwa,Alekhin:2013nda,Ethier:2017zbq}.

Recently, Ji introduced a set of quasi-PDFs for a hadron of momentum $p_z$, say along $z$-direction, and argued that they are equal to corresponding PDFs when the hadron momentum $p_z$ goes to infinity \cite{Ji:2013dva}.
Without setting $p_z\to\infty$, the quasi-PDFs could be factorized to PDFs to all orders in QCD perturbation theory so long as quasi-PDFs can be multiplicatively renormalized \cite{Ma:2014jla}.
Like PDFs, the quasi-PDFs are defined by hadronic matrix elements of two-field correlators with a straight line gauge link between the two fields to ensure the gauge invariance,
\begin{eqnarray}
\label{eq:ftq}
\tilde{f}_{q/p}(\tilde{x},\tilde{\mu}^2,p_z)
&\equiv &
\int \frac{d\xi_z}{2\pi} e^{i \tilde{x}p_z \xi_z}\langle h(p)| \overline{\psi}_q(\xi_z)\,\frac{\gamma_z}{2}
\nonumber\\
&\ & \times
\Phi_{n_z}^{(f)}(\{\xi_z,0\})\, \psi_q(0) | h(p) \rangle
\end{eqnarray}
for quasi-quark distribution with $\xi_0=\xi_\perp=0$, and
\begin{eqnarray}
\label{eq:ftg}
\tilde{f}_{g/p}(\tilde{x},\tilde{\mu}^2,p_z)
&\equiv &
\frac{1}{\tilde{x}p_z}
\int \frac{d\xi_z}{2\pi} e^{i \tilde{x} p_z \xi_z}\langle h(p)| F_{z}^{\ \nu}(\xi_z)\,
\nonumber\\
&\ & \times
\Phi_{n_z}^{(a)}(\{\xi_z,0\})\, F_{z\nu}(0) | h(p) \rangle
\end{eqnarray}
for quasi-gluon distribution with $\nu$ summing over transverse directions.
In Eqs.~(\ref{eq:ftq}) and (\ref{eq:ftg}), $\Phi_{n_z}^{(f,a)}(\{\xi_z,0\}) = \text{exp}[-ig\int_0^{\xi_z} d\eta_z\, A^{(f,a)}_z(\eta_z)]$ are the gauge links with ``$f$'' and ``$a$'' representing fundamental and adjoint representation, respectively.
Unlike PDFs, the two-field correlators of the quasi-PDFs are defined to be off the light-cone and have an equal time separation, which makes it possible to calculate the quasi-PDFs in LQCD \cite{Lin:2014zya,Alexandrou:2015rja,Chen:2016utp,Zhang:2017bzy,Orginos:2017kos}.
However, since the hadron momentum in LQCD calculation is effectively bounded by the lattice spacing, the $p_z\to\infty$ limit is hard to achieve in LQCD calculation, and it is a challenge to control the corrections due to the finite $p_z$ \cite{Ji:2014gla,Ji:2017rah}.  Nevertheless, with the potential to calculate the PDFs from the first principle in QCD by using LQCD and the challenges in doing so, the concept of the quasi-PDFs has generated a lot of interests and activities in both perturbative QCD (PQCD) and LQCD community \cite{Ji:2014hxa,Gamberg:2014zwa,Jia:2015pxx,Li:2016amo,Bacchetta:2016zjm,Radyushkin:2016hsy,Radyushkin:2017gjd,Gbedo:2017eyp,Radyushkin:2017ffo,Carlson:2017gpk,Briceno:2017cpo,Nam:2017gzm,Xiong:2017jtn,Radyushkin:2017cyf,Constantinou:2017sej,Yoon:2017qzo,Rossi:2017muf}.  Besides of the corrections due to the limited range of $p_z$, the key to derive PDFs from the LQCD calculated quasi-PDFs is of two-folds:  (1) being able to renormalize all ultraviolet (UV) divergences of quasi-PDFs nonperturbatively, and (2) ensuring the renormalized quasi-PDFs and PDFs to share the same collinear (CO) divergences.

It was demonstrated in Ref.~\cite{Ma:2014jla} by two of the present authors that quasi-PDFs and corresponding PDFs share the same leading logarithmic CO perturbative divergences to all orders in QCD perturbation theory, if the quasi-PDFs can be multiplicatively renormalized.  With this fact, instead of deriving the PDFs from quasi-PDFs by taking the hadron momentum $p_z\to\infty$, two of the present authors proposed in Ref.~\cite{Ma:2014jla} to extract PDFs by using the PQCD factorization approach from LQCD calculated quasi-PDFs or other LQCD calculable single-hadron matrix elements, so long as these matrix elements could be factorized into the desired PDFs with perturbatively calculable coefficients.  In this PQCD factorization approach, the hard scale is of $\tilde{x}p_z \sim 1/\xi_z \sim \text{GeV}$, and the corrections are power suppressed by the hard scale and characterized by the hadronic matrix elements of high twist operators, just like how experimentally measurable and perturbative factorizable hadronic cross sections are connected to PDFs via PQCD factorization.

Therefore, understanding the renormalizability of the operators defining the quasi-PDFs is the most critically important challenge for extracting PDFs from LQCD calculated quasi-PDFs, reliably.  Since the quasi-PDFs are not defined by twist-2 operators, PDFs and quasi-PDFs have different perturbative UV behavior.  Instead of the logarithmic perturbative UV divergence of PDFs, quasi-PDFs have power UV divergences \cite{Ma:2014jla,Xiong:2013bka}.  Although LQCD calculations of quasi-PDFs are naturally regularized by the lattice spacing $a$, the perturbative power divergence in $1/a$ makes it difficult to take the continuous limit of lattice results to extract the correct PDFs \cite{Ishikawa:2016znu,Chen:2016fxx,Monahan:2016bvm,Alexandrou:2017huk, Chen:2017mzz}.  Under the UV renormalization, it is also possible that the operators defining the quasi-PDFs might mix with other operators, for example, the quark PDFs can mix with gluon PDF.  Therefore, it is very important to find out that not only if the operators defining quasi-PDFs are renormalizable, and but also if the operator mixing under the renormalization, if there is any, could be limited to a closed set of operators.  In this paper, we address the issues concerning the renormalizability of the operators defining quasi-PDFs.


The rest of this paper is organized as following.
In Sec.~\ref{sec:def} we will show that quasi-PDFs defined in Eqs.~\eqref{eq:ftq} and \eqref{eq:ftg} have bad short-distance behavior, while coordinate-space quasi-PDFs are better candidates for extracting PDFs. We also define quasi-PDFs in coordinate-space and explain why its renormalization is difficult. We then study the one-loop expansion of quasi-PDFs in coordinate-space in Sec.~\ref{sec:oneloop}. Using the UV power counting derived in Sec.~\ref{sec:powercounting}, we identify all  perturbative UV divergent regions for the quasi-PDFs.  We find that, once subdivergences are subtracted off, all UV divergences are originated from the integration regions where all loop momenta are large, which is significantly different from the behavior of PDFs.
Based these observations and findings, we prove  in Sec.~\ref{sec:renormalization} that quasi-PDFs in coordinate-space can be multiplicatively renormalized. Most importantly, we found that quasi-PDFs do not mix with other quantities under the running of the renormalization scale, which completes our proof of the renormalizability of coordinate-space quasi-PDFs.  Finally, we give our summary in Sec.~\ref{sec:summary}.
